\subsubsection{Standard Two-by-Two Matrix Realization}
\label{sec:matrix-realization}
The algebra fixed by the three sigma generators and their multiplication
rule is commonly called the real Pauli algebra or the three-dimensional
real Clifford algebra. These names refer here only to the algebra already
defined. Each of its elements can be translated one-to-one into a
two-by-two complex matrix so that both addition and multiplication are
preserved. A one-to-one translation with these properties is called a
faithful matrix representation. The same matrix system also represents the
part of the usual flat space-time algebra built from products containing an
even number of space-time basis directions
\cite{clifford1878,pauli1927,hestenes1966}. Its eight real coefficients are
grouped by the internal pseudoscalar into four complex coefficients: one
scalar and one three-vector. Only the previously defined operations are
used below.

Square brackets and the physical row and column positions are used for
matrices. Thus the generic two-by-two matrix is
\(\smat{a&b\\c&d}\), where every entry is a complex scalar.
Each matrix operation used below is defined separately:
\begin{equation}
\I:=
\begin{bmatrix}
1&0\\
0&1
\end{bmatrix}.
\label{eq:matrix-identity}
\end{equation}
Multiplication by \(\I\) leaves any matrix of the matching size unchanged,
so \(\I\) is the matrix identity.
\begin{equation}
\eqtext{tr}\!\left(
\begin{bmatrix}
a&b\\
c&d
\end{bmatrix}
\right):=a+d.
\label{eq:matrix-trace}
\end{equation}
The sum of the diagonal entries is called the matrix trace.
\begin{equation}
\eqtext{diag}(A,B):=
\begin{bmatrix}
A&0\\
0&B
\end{bmatrix}.
\label{eq:matrix-diag}
\end{equation}
The function \(\eqtext{diag}(\cdot)\) places its inputs on the diagonal and
sets all other entries to zero.
The dimension of \(\I\) is fixed by its surrounding matrix product;
Eq.\ \eqref{eq:matrix-identity} displays the two-by-two case. When more
than two inputs are supplied to \(\eqtext{diag}(\cdot)\), the same rule is
used: the inputs are placed in order on the diagonal and every other entry
is set to zero.
\begin{equation}
\mtrans{\begin{bmatrix}a&b\\c&d\end{bmatrix}}
:=\begin{bmatrix}a&c\\b&d\end{bmatrix}.
\label{eq:matrix-transpose}
\end{equation}
The superscript \(\mathsf T\) exchanges rows and columns and is called the
transpose.
For a scalar entry \(z=x+iy\), the matrix-only conjugate is
\begin{equation}
\mconj{z}:=x-iy.
\label{eq:matrix-entry-conjugate}
\end{equation}
\begin{equation}
\mconj{\begin{bmatrix}a&b\\c&d\end{bmatrix}}
:=\begin{bmatrix}
\mconj{a}&\mconj{b}\\
\mconj{c}&\mconj{d}
\end{bmatrix}.
\label{eq:matrix-conjugate}
\end{equation}
The superscript \(\mathsf C\) changes \(i\) to \(-i\) in every entry and is
called entrywise complex conjugation.
\begin{equation}
\mherm{\begin{bmatrix}a&b\\c&d\end{bmatrix}}
:=\begin{bmatrix}
\mconj{a}&\mconj{c}\\
\mconj{b}&\mconj{d}
\end{bmatrix}.
\label{eq:matrix-hermitian}
\end{equation}
The superscript \(\mathsf{CT}\) applies entrywise conjugation and then
transpose; it is called the Hermitian operation.
\begin{equation}
\adjf{\begin{bmatrix}a&b\\c&d\end{bmatrix}}
:=\begin{bmatrix}d&-b\\-c&a\end{bmatrix}.
\label{eq:matrix-adjugate}
\end{equation}
The displayed exchange and sign change of the entries define the matrix
adjugate.
\begin{equation}
\detf{\begin{bmatrix}a&b\\c&d\end{bmatrix}}:=ad-bc.
\label{eq:matrix-determinant}
\end{equation}
The scalar \(ad-bc\) is the matrix determinant.
Entrywise conjugation and transpose commute, so either order gives the
matrix-only operation \(\mathsf{CT}\). The superscript \(\mathsf C\) is
reserved for matrices and their scalar entries. It is not an additional
operation on the sigma algebra itself. An operation defined by the sigma
rules without choosing a matrix translation is called intrinsic; star and
\(H\) are intrinsic in this sense.

The standard matrix translation is
\begin{equation}
\mat{\sigma}_1=\begin{bmatrix}0&1\\1&0\end{bmatrix},
\quad
\mat{\sigma}_2=\begin{bmatrix}0&-i\\i&0\end{bmatrix},
\quad
\mat{\sigma}_3=\begin{bmatrix}1&0\\0&-1\end{bmatrix}.
\label{eq:pauli-matrices}
\end{equation}
The entry \(i\) is the same central unit defined in
Eq.\ \eqref{eq:basis}; as a scalar matrix entry it obeys \(i^2=-1\).
No second imaginary unit is introduced. The distinction in the matrix
model lies in the operation being applied. Every displayed \(i\), including
the one built into \(\mat{\sigma}_2\), is conjugated by entrywise \(\mathsf C\).
Every sigma generator is additionally flipped by intrinsic star.

The pseudoscalar is represented by the matrix product
\(\mat{\sigma}_1\mat{\sigma}_2\mat{\sigma}_3=i\I\). In this realization,
\begin{equation*}
\mat{\sigma}_n^2=\I,
\qquad
\mat{\sigma}_m\mat{\sigma}_n=-\mat{\sigma}_n\mat{\sigma}_m\quad(m\neq n),
\qquad
\mat{\sigma}_1\mat{\sigma}_2\mat{\sigma}_3=i\I.
\end{equation*}
Thus Eq.\ \eqref{eq:pauli-matrices} is a concrete model of the sigma
generators and their pseudoscalar.

For the generic complex paravector
\(\pv{Z}=S+\vct{V}\cdot\sigv\), its represented matrix is denoted by
\(\mat{M}\):
\begin{equation}
\mat{M}:=
\begin{bmatrix}
S+V_3&V_1-iV_2\\
V_1+iV_2&S-V_3
\end{bmatrix}.
\label{eq:matrix-paravector-map}
\end{equation}
Thus \(\pv{Z}\) is represented by \(\mat{M}\).
Conversely, if \(\mat{M}=\smat{a&b\\c&d}\), then
\begin{equation}
\begin{aligned}
S&=\frac{a+d}{2},&
V_1&=\frac{b+c}{2},\\
V_2&=\frac{c-b}{2i},&
V_3&=\frac{a-d}{2}.
\end{aligned}
\label{eq:matrix-paravector-inverse-map}
\end{equation}
The realization in both directions is therefore given by
Eqs.\ \eqref{eq:matrix-paravector-map} and
\eqref{eq:matrix-paravector-inverse-map}.

For a real paravector
\(\pv{X}=t+\vct{r}\cdot\sigv\in\Pspace(\R)\), with
\(\vct{r}=\vcol{r_1,r_2,r_3}\in\R^3\), the matrix form is
\begin{equation}
\begin{gathered}
\pv{X}=t+\vct{r}\cdot\sigv\\
\eqtext{is represented by}\\
\begin{bmatrix}
t+r_3&r_1-ir_2\\
r_1+ir_2&t-r_3
\end{bmatrix},\\
\detf{\pv{X}}=t^2-\vct{r}^2.
\end{gathered}
\label{eq:matrix-X}
\end{equation}
The real space-time expression \(t^2-\vct{r}^2\), conventionally called the
Minkowski quadratic form and defined in
Eq.\ \eqref{eq:real-paravector-metric}, is reproduced by
Eq.\ \eqref{eq:matrix-X}, together with
\begin{align*}
\eqtext{tr}(\mat{M})&=2t,\\
x^2-\eqtext{tr}(\mat{M})x+\detf{\mat{M}}&=0,\\
x&=t\pm\lVert\vct{r}\rVert.
\end{align*}
Here \(\mat{M}\) denotes the displayed matrix representing \(\pv{X}\), and
\(x\) is only the polynomial input. The last two lines show that
\(t+\lVert\vct{r}\rVert\) and \(t-\lVert\vct{r}\rVert\) are the two roots
of the displayed matrix equation. For each root there is a nonzero column that is
merely multiplied by that root when acted on by the matrix. Scalars with
this property are called eigenvalues. They give the matrix form of the
split into scalar part \(t\) and vector magnitude
\(\lVert\vct{r}\rVert\).

The sigma-basis action not captured by entrywise \(\mathsf C\) is accounted for by
the fixed matrix
\begin{equation}
\Tmat:=
\begin{bmatrix}
0&1\\
-1&0
\end{bmatrix}.
\label{eq:matrix-T}
\end{equation}
It satisfies \(\Tmat^{-1}=-\Tmat\).
The complete intrinsic-to-matrix dictionary is
\begin{equation}
\begin{aligned}
\pv{Z}
&\quad\eqtext{is represented by}\quad \mat{M},\\
\conj{\pv{Z}}
&\quad\eqtext{is represented by}\quad
\conj{\mat{M}},\\
\conj{\mat{M}}
&=\Tmat\,\mconj{\mat{M}}\,\Tmat^{-1}
\\
&=\begin{bmatrix}
\mconj{d}&-\mconj{c}\\
-\mconj{b}&\mconj{a}
\end{bmatrix},\\
\herm{\pv{Z}}
&\quad\eqtext{is represented by}\quad
\herm{\mat{M}},\\
\herm{\mat{M}}
&=\mherm{\mat{M}}
\\
&=\begin{bmatrix}
\mconj{a}&\mconj{c}\\
\mconj{b}&\mconj{d}
\end{bmatrix},\\
\adjf{\pv{Z}}
&=\cherm{\pv{Z}},\\
\adjf{\pv{Z}}
&\quad\eqtext{is represented by}\quad
\adjf{\mat{M}},\\
\adjf{\mat{M}}
&=\Tmat\,\mtrans{\mat{M}}\,\Tmat^{-1}
\\
&=\begin{bmatrix}d&-b\\-c&a\end{bmatrix},\\
\detf{\pv{Z}}
&\quad\eqtext{is represented by}\quad
\detf{\mat{M}}=ad-bc,\\
0&\quad\eqtext{is represented by}\quad
\begin{bmatrix}0&0\\0&0\end{bmatrix},\\
1&\quad\eqtext{is represented by}\quad\I,\\
\pv{Z}^{-1}
&\quad\eqtext{is represented by}\quad
\mat{M}^{-1},\\
\mat{M}^{-1}
&=\frac{\adjf{\mat{M}}}{\detf{\mat{M}}},
\qquad \detf{\mat{M}}\neq0.
\end{aligned}
\label{eq:matrix-operation-dictionary}
\end{equation}
Thus \(\herm{\mat{M}}=\mherm{\mat{M}}\) is the ordinary Hermitian
adjoint. A matrix is Hermitian precisely when
\(\herm{\mat{M}}=\mat{M}\), while the matrix adjugate is always denoted by
\(\adjf{\mat{M}}\). Crucially, intrinsic star is not matrix \(\mathsf C\):
\(\conj{\mat{M}}\neq\mconj{\mat{M}}\) in general.

The star on the left is the intrinsic sigma conjugate transported into
the matrix realization; the \(\mathsf C\) on the right is only entrywise complex
conjugation. They agree on central complex scalars, but not on a generic
paravector. When \(\pv{Z}\) is represented by \(\mat{M}\), the direct bridge
is that \(\herm{\pv{Z}}\) is represented by \(\mherm{\mat{M}}\). Intrinsic star
instead requires the basis correction
\(\conj{\mat{M}}=\Tmat\mconj{\mat{M}}\Tmat^{-1}\).

The extraction functions have the equally direct matrix form
\begin{equation}
\begin{aligned}
\eqtext{scalar}(\pv{Z})
&\quad\eqtext{is represented by}\quad\tfrac12\eqtext{tr}(\mat{M}),\\
\eqtext{vector}(\pv{Z})\cdot\sigv
&\quad\eqtext{is represented by}\quad
\mat{M}-\tfrac12\eqtext{tr}(\mat{M})\I,\\
\eqtext{real}(\pv{Z})
&\quad\eqtext{is represented by}\quad
\tfrac12\bigl(\mat{M}+\herm{\mat{M}}\bigr),\\
\eqtext{img}(\pv{Z})
&\quad\eqtext{is represented by}\quad
\frac{\mat{M}-\herm{\mat{M}}}{2i}.
\end{aligned}
\label{eq:matrix-extraction-dictionary}
\end{equation}
Consequently, the algebraic determinant is exactly the
ordinary two-by-two matrix determinant:
\begin{equation}
\mat{M}\adjf{\mat{M}}
=\adjf{\mat{M}}\mat{M}
=\detf{\mat{M}}\I.
\label{eq:matrix-adjugate-product}
\end{equation}

The operations are distinguished by whether they conjugate complex scalar
coefficients and whether they reverse the order of a product. Matrix
\(\mathsf C\) conjugates the coefficients and preserves product order.
Matrix \(\mathsf T\) leaves the coefficients unchanged and reverses product
order. Their composition \(\mathsf{CT}\), represented in the sigma algebra
by \(H\), both conjugates coefficients and reverses product order. Intrinsic
star also conjugates coefficients but preserves product order. The combined
intrinsic operation \(*H\), represented by the matrix adjugate, leaves
coefficients unchanged and reverses product order. Consequently,
Eq.\ \eqref{eq:hermitian-adjugate-identity} is represented by
\begin{equation*}
\herm{\mat{M}}
=\conj{\adjf{\mat{M}}}
=\adjf{\conj{\mat{M}}},
\end{equation*}
where the star in \(\conj{\mat{M}}\) remains the induced sigma conjugate,
not entrywise \(\mathsf C\).

The distinction is made concrete by the basis actions:
\begin{equation}
\begin{aligned}
\mconj{\mat{\sigma}_1}&=\mat{\sigma}_1,
&
\mconj{\mat{\sigma}_2}&=-\mat{\sigma}_2,
&
\mconj{\mat{\sigma}_3}&=\mat{\sigma}_3,\\
\mtrans{\mat{\sigma}_1}&=\mat{\sigma}_1,
&
\mtrans{\mat{\sigma}_2}&=-\mat{\sigma}_2,
&
\mtrans{\mat{\sigma}_3}&=\mat{\sigma}_3,\\
\mherm{\mat{\sigma}_n}&=\mat{\sigma}_n,
&
\conj{\mat{\sigma}_n}&=-\mat{\sigma}_n,
&
\cherm{\mat{\sigma}_n}&=-\mat{\sigma}_n.
\end{aligned}
\label{eq:matrix-basis-actions}
\end{equation}
On these three individual generator matrices, the same sign table is obtained
from \(\mathsf C\) and \(\mathsf T\). Distinct operations are nevertheless represented:
\(\mconj{(\mat{A}\mat{B})}=\mconj{\mat{A}}\mconj{\mat{B}}\), whereas
\(\mtrans{(\mat{A}\mat{B})}=\mtrans{\mat{B}}\mtrans{\mat{A}}\). Only the represented second
generator is changed by entrywise \(\mathsf C\), while all three are changed by
intrinsic star. Every \(\mat{\sigma}_n\) is fixed by their composition
\(\mathsf{CT}\), exactly matching intrinsic \(H\). Thus \(\mathsf C\) depends
on the chosen matrix entries
and is distinct from both sigma-algebra operations. In particular,
it is not intrinsic star. For \(n\neq m\),
\[
\herm{\mat{\sigma}_n\mat{\sigma}_m}
=\mat{\sigma}_m\mat{\sigma}_n
=-\mat{\sigma}_n\mat{\sigma}_m,
\]
exactly as required by the intrinsic order-reversal rule.

\subsubsection{Auxiliary Block Map}
A scalar-coefficient paravector \(\pv{Z}\in\Pspace\) is packaged
into the \(2\times2\) block map
\begin{equation}
\label{eq:W-map}
\Wmat(\pv{Z}):=
\begin{bmatrix}
0&\pv{Z}\\
\adjf{\pv{Z}}&0
\end{bmatrix}.
\end{equation}
\begin{equation}
\Wmat(\pv{Z})^2
=\eqtext{diag}(\detf{\pv{Z}},\detf{\pv{Z}})
=\detf{\pv{Z}}\I.
\label{eq:W-square}
\end{equation}
The block map \(\Wmat(\pv{Z})\) is a \(2\times2\) matrix over the paravector
algebra. After each block is replaced by its Pauli matrix, it is a
\(4\times4\) complex matrix. Addition and multiplication by ordinary
scalars are preserved, so the packaging map is linear. General products are
not preserved, so it is not an additional representation of the algebra.
Equation \eqref{eq:W-square} is obtained because its
diagonal block products are reduced to the determinant in
Eq.\ \eqref{eq:minkowski-form}; the off-diagonal blocks are zero
\cite{hestenes1966,doran2003}.

If \(\detf{\pv{Z}}\neq0\), Eq.\ \eqref{eq:W-square} shows that the powers of
the block matrix are governed by the scalar equation
\(x^2-\detf{\pv{Z}}=0\). Its two roots are distinct; selecting one square
root fixes their labels. An invertible change of basis
can then put the block matrix into diagonal form; this property is called
diagonalizability. If
\(\detf{\pv{Z}}=0\) and \(\pv{Z}\neq0\), then
\(\Wmat(\pv{Z})^2=0\) while \(\Wmat(\pv{Z})\neq0\). A nonzero matrix whose
positive power is zero is called nilpotent; such a matrix cannot be put
into diagonal form. For blocks containing differential operations or other
quantities whose order cannot be exchanged, the unreduced statement is
\(\Wmat(\pv{Z})^2
=\eqtext{diag}(\pv{Z}\cherm{\pv{Z}},\cherm{\pv{Z}}\pv{Z})\).
A separate operator calculation is required before either ordered product can
be called a determinant or both can be reduced to one scalar multiple of the
identity.
